\section{Review Exercises}

\subsection{Elementary computational and construction exercises}

\begin{enumerate}
\item Draw two perpendicular axes, choose an origin and a common unit, and plot
\[
A=(3,2),\qquad B=(-2,4),\qquad C=(-3,-1),\qquad D=(2,-3).
\]
For each point, state its horizontal and vertical signed displacement from the origin.

\item For the point $P=(-4,3)$:
\begin{enumerate}
\item draw its orthogonal projections onto both axes;
\item write the coordinates of the two projection points;
\item compute the distance from $P$ to each axis.
\end{enumerate}

\item Compute the distance between each pair of points:
\[
A=(1,2),\ B=(4,6),
\qquad
C=(-2,1),\ D=(4,1).
\]
Draw the corresponding right triangles whose legs produce the distance calculation.

\item Find all points on the $x$-axis whose distance from $P=(2,3)$ is $5$. Verify each answer with the distance formula.

\item A rectangle has vertices
\[
A=(-2,-1),\qquad B=(3,-1),\qquad C=(3,2).
\]
Find the fourth vertex, compute the side lengths, and calculate the length of both diagonals.

\item Reflect the point $P=(3,-4)$:
\begin{enumerate}
\item across the $x$-axis;
\item across the $y$-axis;
\item across the origin.
\end{enumerate}
Plot all four points and record the coordinate changes.

\item The coordinate unit is changed so that one new unit equals two old units. A fixed point originally has coordinates $(6,-4)$. Find its coordinates in the new scale. Repeat when one new unit equals one half of an old unit.

\item Construct with straightedge and compass:
\begin{enumerate}
\item the perpendicular bisector of a segment $AB$;
\item a line through a point $P$ perpendicular to a given line;
\item a line through $P$ parallel to the given line, using the perpendicular construction twice.
\end{enumerate}

\item In a right triangle the legs have lengths $6$ and $8$.
\begin{enumerate}
\item compute the hypotenuse;
\item compute the sine and cosine of each acute angle from the side lengths;
\item reproduce the same triangle in a Cartesian system with one vertex at the origin.
\end{enumerate}

\item The points $A=(-1,2)$ and $B=(5,2)$ lie on a horizontal line. Divide the segment $AB$ into six equal parts using the coordinate scale, and list the coordinates of all division points.
\end{enumerate}

\subsection{Conceptual and reasoning exercises}

\begin{enumerate}
\item Explain why points, lines, circles, perpendicularity, and congruence are meaningful before coordinates are introduced. Which new capability is added by coordinates?

\item A student says: ``The point $P$ is the ordered pair $(2,3)$.'' Correct this statement by distinguishing the geometric object from its coordinate description.

\item Two students draw different Cartesian systems on the same plane and assign different coordinate pairs to the same point. Explain why neither description is necessarily wrong. List the choices that may differ between their systems.

\item Explain why the construction of a right angle does not require a numerical measure such as $90^\circ$. What does numerical angle measure add to the geometric construction?

\item Starting from the Pythagorean theorem, justify the distance formula
\[
d(P,Q)=\sqrt{(x_2-x_1)^2+(y_2-y_1)^2}.
\]
Identify the geometric meaning of the two coordinate differences.

\item A point has coordinates $(4,3)$ in one system. The common unit is doubled, but the axes and origin are unchanged. Explain why its coordinates change while its geometric position and distance from the origin do not.

\item Why must the two Cartesian axes use a common unit if the usual distance formula is to represent Euclidean length directly? Describe what would go wrong if the horizontal and vertical scales were different.

\item Compare the following three statements and give one example of each from the chapter:
\begin{enumerate}
\item a Euclidean postulate;
\item a derived construction;
\item a coordinate formula.
\end{enumerate}
Explain why they have different logical roles.

\item A drawing appears to show that two lines are perpendicular. Explain why visual appearance alone is not a proof. Give one construction-based and one coordinate-based way to justify perpendicularity.

\item Give a coherent account of the chapter's central transition: begin with the geometric objects already available, describe the choices that create a Cartesian system, and explain how arithmetic becomes a language for geometry without replacing the geometry itself.
\end{enumerate}
